%=============================================================================
% RESULTS
%=============================================================================
\section{Results}
\label{sec:results}

All numbers are RMSE in weekly cases on the corrected dataset, from the 810 runs
of the reproduction notebook. Validation decides and test is reported beside it.
Unless stated otherwise, pairs are on the (origin, seed) unit of the frozen three
origins.

\begin{figure*}[t]
  \centering
  \includegraphics[width=\linewidth]{figures/fig_encoders}
  \caption{Six encoders behind three heads (Table~\ref{tab:encoders}). The direct
  head separates encoders that work from encoders that fail; the SEIR decoder is
  worse than persistence on all of them; the gated SEIR head puts all six at the
  same level, below persistence (dashed) on validation and test.}
  \label{fig:encoders}
\end{figure*}

\subsection{Six encoders, one harness}
\label{sec:encoders}
Table~\ref{tab:encoders} and Figure~\ref{fig:encoders} give the comparison the
SEIR-GNN question needs, with every encoder sharing folds, loss, early stopping
and metric.

\textbf{Architecture matters little, and the graph less.} Behind the direct head
AAGCN, ASTGCN and the LSTM score 16.71, 16.83 and 16.91: a spread of 0.20, and
all three beat persistence (17.86). The LSTM has no message passing at all. A
GCN beats the same network without message passing by only 0.06 (3/9,
Table~\ref{tab:levers}). STGAT, A3TGCN and DCRNN score 22.93, 28.63 and 34.64,
all worse than persistence. A linear model fitted to their training residuals
recovers 48\% and 60\% of STGAT's and A3TGCN's validation residual variance, but
only 3--7\% of the working encoders'. The failing encoders are underfitting,
limited by their wiring rather than by the data.

\textbf{The SEIR decoder fails on every encoder.} Routing the forecast entirely
through the simulator scores 24.01--25.75, worse than persistence on all six and
7.3--7.4 worse than the direct head on each working encoder. The deficit is
structural, and we measured it without training by inverting the simulator for
the force of infection that reproduces each observed count:
(i) with transmission switched off, the $E$ already seeded from recent cases
still emits \emph{more} cases than observed for 14--16\% of training targets,
which are therefore unreachable at any network output (39--41\% of cells need a
force of infection of zero);
(ii) the required force of infection is only $r^2=0.22$--$0.26$ predictable from
case history, against 0.81--0.82 for the direct target;
(iii) it is not susceptible depletion, since only 0.7--4.9\% of cells reach the
susceptible clamp.

\textbf{The gated SEIR head repairs the failing encoders.} Letting the
simulator supply only the departure from persistence brings all six encoders to
17.43--17.69, below persistence on validation (17.86) and on test (35.50--35.84
against 36.02). For the three failing encoders this is a large repair, on 9/9
paired runs each: STGAT $-5.30$ (test $-26.16$), A3TGCN $-10.95$ ($-24.90$) and
DCRNN $-16.97$ ($-38.01$). It costs the working encoders 0.68--0.78 on
validation. Read at face value, this is the physics-informed result the study
set out to find. The next section shows that it is not.

\subsection{Anchor, gate or physics?}
\label{sec:rescue}
The gated SEIR head differs from the direct head in three ways at once: it
forecasts a departure from last week's value (the anchor), it shrinks that
departure through a learned gate, and the departure comes from the SEIR
simulator. Only the third is physics. A pre-registered control separates them
by training all six encoders behind four heads that add these one at a time:
direct, residual (anchor), gated (anchor and gate) and gated SEIR (all three).
The rescue is credited to physics only if gated SEIR beats its no-physics twin on
at least 7 of 9 runs with $\padj<0.05$ on each failing encoder. Our stated
prediction was that the anchor and gate carry most of the rescue.

Table~\ref{tab:rescue} gives the result: \textbf{the rescue is the anchor, not
the physics}. Anchor and gate alone account for 96--112\% of the validation
repair on each failing encoder, and a plain residual head, with no gate and no
simulator, already brings them to 17.47--17.85. Adding the simulator to the
gated head earns credit on none of the three: it is significantly \emph{worse} on
STGAT ($+0.66$, 1/9, $\padj=0.023$), ties on A3TGCN ($-0.00$, 6/9), and helps
DCRNN by $-0.68$ on only 3/9 runs ($\padj=0.38$). On the three working encoders
the simulator costs 0.75--0.86 on 0 or 1 of 9 runs. The best head for AAGCN is
the gated head without physics (16.57 against the direct head's 16.71), a
comparison we did not pre-register and report only as a lead.

On test the picture reverses: the SEIR head beats its no-physics twin on the
failing encoders (STGAT $-0.62$, 8/9; A3TGCN $-1.66$ and DCRNN $-3.50$, both
9/9). That is the recurring disagreement of \S\ref{sec:valtest}. It was not used
to decide anything, and we do not claim it.

\subsection{What moves the metric}
\label{sec:levers}
Table~\ref{tab:levers} pairs every lever we tried against its own control, and
Figure~\ref{fig:levers} shows the same deltas. The \textbf{negative-binomial
likelihood} is the largest controlled gain: $-0.76$, 9/9, $\padj=0.014$, larger
than any architecture difference. It corrects a bias, since squared error on
$\log(1+y)$ targets a median while RMSE scores a mean. \textbf{Seasonal features}
give a similar mean change ($-0.76$) but win only 6/9. Combined, the best model
$B$ (AAGCN, direct head, NB, seasonal features) scores 15.51 against persistence's
17.86: $-2.35$ on 9/9 runs.

That is the headline a benchmark would report. It does not hold out of sample: on
test, $B$ scores 37.54 against persistence's 36.02, and the NB likelihood is
\emph{worse} on test by 1.61. We return to this in \S\ref{sec:valtest}.

\begin{figure}[t]
  \centering
  \includegraphics[width=\linewidth]{figures/fig_levers}
  \caption{Every lever's validation delta against its own control
  (Table~\ref{tab:levers}); win counts on the right. Green: negative and
  $\ge$7/9 wins. Grey: negative but inconsistent. Red: worse.}
  \label{fig:levers}
\end{figure}

\subsection{SEIR-GNN against SEIR-LSTM}
\label{sec:seirgnn}
The comparison the research question poses is a graph encoder against the LSTM,
behind the same gated SEIR head (Table~\ref{tab:nine}). On three origins the
graph wins on validation: ASTGCN $-0.54$ (9/9) and AAGCN $-0.14$ (8/9,
$\padj=0.020$ in Table~\ref{tab:levers}). Behind the direct head the same graph
gains less and less consistently ($-0.27$, 6/9) and loses on test ($+1.96$, 1/9).
So what the graph adds appears specific to estimating a transmission quantity.

The nine disjoint origins do not support a general claim. ASTGCN's advantage
holds on validation ($-0.62$, 8/9, $p=0.012$) and vanishes on test ($+0.07$,
6/9); AAGCN's reverses the other way ($+0.18$ on validation, $-0.30$ on test,
5/9, $p=0.078$). Neither is significant on the metric where the other is. We
therefore do not claim that the SEIR-GNN beats the SEIR-LSTM: it does so
consistently on three origins on validation, and on nine origins it is not robust
to the encoder or the metric.

\begin{figure*}[t]
  \centering
  \includegraphics[width=\linewidth]{figures/fig_nine}
  \caption{Nine disjoint origins, per-origin RMSE on a log scale. Origin 0.40's
  validation span contains the 2017 epidemic and dominates every validation
  mean; read win counts, not means.}
  \label{fig:nine}
\end{figure*}

\subsection{Physics as structure: penalty and metapopulation coupling}
\label{sec:physicsstruct}
Two further physics-informed structures were tested against $B$. The
\textbf{spatial physics penalty} was the one significant physics gain in our
earlier work on the benchmark array ($-0.044$ on STGAT, 57/60 runs, reported in
the Phase-2 study). On $B$ it keeps its sign ($-0.06$) but wins only 6/9; its
log form is worse ($+0.08$, 2/9). The \textbf{metapopulation SEIR head}, which
couples districts inside the dynamics rather than through the initial state, is
\emph{worse} than the gated head ($+0.20$, 1/9, $\padj=0.020$), and on nine
origins it ties the SEIR-LSTM on test ($-0.01$, 4/9). The reason is in the data:
on the district graph, Moran's $I$ of weekly log growth is $-0.03$, against 0.29
for log cases. Levels are spatially clustered; transmission changes are not, so
a learned coupling has nothing to encode.

\subsection{The ceiling is informational}
\label{sec:ceiling}
Table~\ref{tab:levers} records what did not work. No architecture change
improves validation by more than 0.15 (a learned adjacency, 5/9), and several
make it worse: the nonlinear head by 0.28, the residual head with NB by 0.47 and
all changes combined by 0.55. The literature's remedies for weak
spatio-temporal models are neutral or worse: RevIN costs 1.47 (0/9) and
$k$-nearest-neighbour analogues with no network cost 1.01 (0/9). Data-side levers
fail too. Training on a quarter of the windows costs only 0.25, and on
three-quarters 0.11, so the model is not short of data. TimeGAN augmentation
costs 0.60 (0/9) and truncates the tail: its largest generated week is 2,036
cases against a real 2,631. Climate costs 0.18, 0.19 and 0.47 at lags 2--4,
2--13 and 2--25. Gradient-boosted trees do gain from climate ($-0.54$, 6/9), but
remain far behind $B$ (17.40).

These failures share a cause. The validation residuals of the working encoders,
the SEIR-LSTM and $k$-NN, which has no network and no training, correlate
0.91--0.97 with $B$'s: the models make the same errors. The remaining error is a
lag on moves. $B$ under-predicts rising weeks by 11.5 cases and over-predicts
falling weeks by 10.1, and the top 5\% of district-weeks carry 54--67\% of every
model's squared error. That is what a mean-seeking objective produces on a
near-random-walk target whose next move is not in the inputs.

\subsection{Validation and test disagree, systematically}
\label{sec:valtest}
Across the eight families of runs, the configuration validation selects is worse
than persistence on test in five, and the best configuration on test is anchored
to the last observation (gated SEIR, residual, gated or $k$-NN) in six.
Examples: the residual head with NB (34.59 against $B$'s 37.54), $k$-NN (34.34),
the SEIR-LSTM (34.97), and on nine origins the gated SEIR-GNN at 31.17 against
$B$ at 34.00 and persistence at 31.76. Nothing here was selected on test, and on
nine origins the gated SEIR-GNN's test advantage over persistence is not
significant ($-0.59$, 6/9, $p=0.68$). A pattern across most families is
nonetheless evidence about the selection protocol. A 30-week validation span
immediately before each origin appears to reward fitting that span's level,
which does not carry into the next period. Acting on it needs a new
pre-registered protocol and a test set nothing here has touched, which only data
after March 2024 can provide.
